\chapter{Conclusiones del informe}

\section{Evolución y democratización de las amenazas digitales}
Tras analizar detalladamente los diversos casos de estudio presentados a lo largo de este informe, resulta evidente que el panorama de la ciberseguridad ha experimentado una transformación drástica en los últimos años. Las amenazas han dejado de ser incidentes aislados protagonizados por atacantes solitarios para convertirse en operaciones a escala industrial, orquestadas por verdaderos sindicatos del crimen organizado y actores patrocinados por Estados. La convergencia de tecnologías emergentes, como la Inteligencia Artificial, con vectores de ataque clásicos ha provocado que la complejidad estructural de los incidentes escale de manera exponencial. 

El modelo de "Cibercrimen como Servicio" (CaaS) ha democratizado el acceso a herramientas destructivas. Operaciones como las gigantescas botnets del Internet de las Cosas (capaces de generar tráfico DDoS de hasta 30 Tbps) demuestran que ya no es necesario poseer altos conocimientos técnicos para paralizar las infraestructuras de un país; basta con tener los recursos financieros para alquilar estos servicios en foros clandestinos.

\section{Análisis transversal de los vectores de ataque}

\subsection{La Inteligencia Artificial como catalizador del engaño}
Uno de los descubrimientos más alarmantes en el estudio del phishing moderno es la integración profunda de la IA generativa. El caso en el que se defraudaron más de 19 millones de euros utilizando \textit{deepfakes} demuestra que el clásico correo mal redactado ha quedado en el pasado. Los atacantes ahora son capaces de crear campañas a largo plazo (como la estafa \textit{Pig Butchering}), donde suplantan a figuras públicas con un nivel de realismo que resulta indistinguible para el ciudadano medio. Esto nos lleva a concluir que las defensas puramente tecnológicas son insuficientes; es necesaria una reeducación masiva de la sociedad para fomentar un escepticismo digital crítico frente a la hiperrealidad de los contenidos en la red.

\subsection{El secuestro digital: Del archivo al hardware}
La evolución del software malicioso ha ramificado sus estrategias de monetización. Por un lado, tenemos el impacto tradicional y devastador del ransomware en grandes corporaciones, evidenciado en el ataque a DaVita. Este incidente confirma que los sectores críticos (como el sanitario) son presas favoritas por su baja tolerancia a la inactividad operativa y la sensibilidad de sus datos (PHI), lo que empuja a la doble extorsión.

Por otro lado, se ha observado una mutación espeluznante en el entorno móvil con amenazas como SpyLoan. Aquí, el ransomware se transforma en un "secuestro de hardware" disfrazado de solución financiera. El atacante ya no cifra los archivos, sino que bloquea completamente el terminal, utilizando simultáneamente los datos personales para coaccionar y humillar a la víctima en su entorno social. Este enfoque pone de manifiesto que las vulnerabilidades socioeconómicas están siendo explotadas directamente mediante la tecnología.

\section{La fragilidad del software y la amenaza geopolítica}
El estudio de las inyecciones SQL y el Cross-Site Scripting (XSS) en este informe revela una tendencia crítica: el ecosistema del software actual, incluso aquel considerado maduro y ampliamente confiable (como PostgreSQL o Zimbra), oculta debilidades que son aprovechadas activamente en conflictos geopolíticos.

El ciberespionaje ha encontrado en estas vulnerabilidades su mejor aliado. Casos como el del grupo "Silk Typhoon" atacando al Departamento del Tesoro de EE. UU. a través de un \textit{zero-day} en una librería de bases de datos, o el grupo "APT28" inyectando \textit{payloads} silenciosos en correos electrónicos de instituciones ucranianas, demuestran que ninguna capa del software está a salvo. Una inyección SQL no neutralizada ya no solo pone en riesgo una base de datos local, sino que compromete la seguridad nacional de potencias mundiales mediante la ejecución remota de código (RCE). Esto subraya la urgencia de establecer auditorías de código abierto mucho más rigurosas y la adopción universal de arquitecturas de "Confianza Cero" (Zero Trust).

\section{\textquestiondown Qué nos depara el futuro de la ciberseguridad?}
A partir de la redacción de este documento, podemos extraer varias lecciones y previsiones para el futuro a corto y medio plazo en el sector de la Seguridad de la Información:

\begin{enumerate}
    \item \textbf{La omnipresencia del IoT como arma de doble filo:} Mientras los fabricantes prioricen el coste y la facilidad de uso sobre la seguridad desde el diseño ("Security by Design"), de dispositivos domésticos seguirán siendo asimilados en botnets globales. Se requerirá legislación internacional estricta que prohíba contraseñas por defecto y obligue a un ciclo de vida de parches.
    \item \textbf{Colaboración policial internacional:} Casos de éxito como la "Operación Secure" de Interpol confirman que la única forma de desmantelar complejas infraestructuras de \textit{infostealers} (que facilitan intrusiones mayores) es la cooperación transfronteriza y la colaboración directa con empresas privadas de ciberseguridad.
    \item \textbf{La IA defensiva como respuesta imperativa:} Del mismo modo que los atacantes usan IA para automatizar campañas e inyectar código polimórfico, los equipos del \textit{Blue Team} deben delegar sus operaciones de detección de intrusos (IDS), análisis heurístico y triaje de incidentes a inteligencias artificiales capaces de operar a la velocidad de la máquina, neutralizando, por ejemplo, los \textit{deepfakes} antes de que lleguen a los clientes de correo.
\end{enumerate}

\section{Valoración final}
En conclusión, la práctica 8 ha puesto de relieve que la seguridad digital ya no puede ser tratada como un departamento aislado y puramente técnico dentro de las organizaciones, ni como un aspecto secundario en la vida diaria de las personas. La interconexión total del mundo moderno ha provocado que cualquier vulnerabilidad técnica (CWE) escale rápidamente hacia un ataque crítico que impacta la economía de familias, paraliza hospitales y vulnera la soberanía de los Estados. 

Para poder hacer frente a las ciberamenazas de esta década y más allá, es imprescindible adoptar un enfoque holístico que integre el desarrollo seguro de software y unos mecanismos de parcheo veloces e inmediatos. Solo mediante la sinergia entre tecnología, educación, leyes robustas e inversión proactiva podremos construir un ciberespacio más resiliente.